\chapter{A1: Grundlagen - Wortstellung}
\label{cha:A1Wortstellung}

\section{Einfache Wortstellung}
\label{sec:EinfacheWortstellung}

Im Deutschen gilt die Grundregel: Subjekt - Verb - Objekt.

\bigskip
\noindent\textbf{Wichtig:} Die Wortstellung im einfachen Hauptsatz folgt einem festen Muster.

\begin{center}
\begin{tabular}{|c|c|c|c|c|}
\hline
Position 1 & Position 2 & Position 3 & Position 4 & Position 5+ \\ \hline
Subjekt & finites Verb & Objekt & Adverb & ... \\ \hline
Ich & esse & Brot & heute. & \\ \hline
Er & liest & ein Buch. & & \\ \hline
Wir & lernen & Deutsch & schnell. & \\ \hline
\end{tabular}
\end{center}

\bigskip
\noindent\textbf{Tipp:} Das finite Verb (die konjugierte Form) steht immer an Position 2!

\chapterend

\chapter{A1: Präsens der regelmäßigen Verben}
\label{cha:A1Praesens}

\section{Konjugation der regelmäßigen Verben}

Die regelmäßigen (schwachen) Verben folgen einem festen Muster:

\begin{center}
\begin{tabular}{|l|c|c|c|c|c|}
\hline
Person & Endung & machen & sagen & arbeiten & kaufen \\ \hline
ich & -e & mache & sage & arbeite & kaufe \\ \hline
du & -st & machst & sagst & arbeitest & kaufst \\ \hline
er/sie/es & -t & macht & sagt & arbeitet & kauft \\ \hline
wir & -en & machen & sagen & arbeiten & kaufen \\ \hline
ihr & -t & macht & sagt & arbeitet & kauft \\ \hline
sie/Sie & -en & machen & sagen & arbeiten & kaufen \\ \hline
\end{tabular}
\end{center}

\bigskip
\noindent\textbf{A1-Level:} Diese Tabelle ist für A1-Niveau grundlegend!

\chapterend

\chapter{A1: Artikel und Nomen}
\label{cha:A1ArtikelNomen}

\section{Die drei Geschlechter}

Im Deutschen hat jedes Nomen ein Geschlecht (Genus):

\begin{center}
\begin{tabular}{|c|c|c|}
\hline
Maskulin (m) & Feminin (f) & Neutral (n) \\ \hline
der & die & das \\ \hline
der Mann & die Frau & das Kind \\ \hline
der Tisch & die Lampe & das Haus \\ \hline
der Hund & die Katze & das Auto \\ \hline
der Stuhl & die Tür & das Fenster \\ \hline
\end{tabular}
\end{center}

\bigskip
\noindent\textbf{Achtung:} Das Geschlecht muss man lernen!

\chapterend

\chapter{A1: Perfekt}
\label{cha:A1Perfekt}

\section{Bildung des Perfekt}

Perfekt = haben/sein (Konjugiert) + Partizip II

\bigskip
\noindent\textbf{Regel:}
\begin{itemize}
    \item Die meisten Verben: haben + Partizip II
    \item Verben der Bewegung: sein + Partizip II
\end{itemize}

\section{Partizip II der regelmäßigen Verben}

\begin{center}
\begin{tabular}{|l|c|c|}
\hline
Infinitiv & Partizip II & Regel \\ \hline
machen & gemacht & ge- + Stamm + -t \\ \hline
sagen & gesagt & ge- + Stamm + -t \\ \hline
arbeiten & gearbeitet & ge- + Stamm + -et \\ \hline
kaufen & gekauft & ge- + Stamm + -t \\ \hline
spielen & gespielt & ge- + Stamm + -t \\ \hline
\end{tabular}
\end{center}

\section{Perfekt mit sein}

Diese Verben bilden Perfekt mit sein:

\begin{center}
\begin{tabular}{|l|c|c|}
\hline
Infinitiv & Perfekt & Bedeutung \\ \hline
gehen & ist gegangen & went \\ \hline
kommen & ist gekommen & came \\ \hline
fahren & ist gefahren & drove \\ \hline
fliegen & ist geflogen & flew \\ \hline
laufen & ist gelaufen & ran \\ \hline
schwimmen & ist geschwommen & swam \\ \hline
bleiben & ist geblieben & stayed \\ \hline
\end{tabular}
\end{center}

\chapterend